# Maia Link CLI

---
title: maia link
description: Link a local skill folder into a coding agent so it loads your live copy
weight: 70
---

`maia link` creates a link from your local skill folder into the skills folder a coding agent reads (`~/.claude/skills/` or `~/.agents/skills/`).
The agent then loads the skill straight from your folder, so any edits you make show up right away without re-installing.

This is the main step when you're building a skill: link your folder, use the skill in Claude Code, Codex, or OpenCode, edit it, and repeat.
When you're done, remove the link with [`maia unlink`](../unlink/), and check on it any time with [`maia status`](../status/).

## Usage

```bash
maia link [-h] [--harness HARNESS ...] [path]
```

If you don't give a `path`, it uses the current folder.

## Arguments

| Argument | Default | Notes |
| --- | --- | --- |
| `path` | current folder | Path to the local skill folder to link. The skill is named after this folder (see [Name](#name)). |

## Options

| Flag | Default | Notes |
| --- | --- | --- |
| `--harness HARNESS` | every harness | One or more coding agents to link into, space-separated. Omit it to link into all of them. |
| `-h`, `--help` | none | Print help and exit. |

### `--harness`

`--harness` selects which coding agent(s) the link is created in.
It is variadic, so you can pass several at once, space-separated.
Omit it and the link is created in every harness at once.

| Value | Skills folder |
| --- | --- |
| `claude` | `~/.claude/skills/` |
| `codex` | `~/.agents/skills/` |
| `opencode` | `~/.claude/skills/` (OpenCode reads the same layout as Claude Code) |

Linking is all-or-nothing across the harnesses you name.
If the link can't be created in one of them, it isn't created in any of them (see [When the Name Is Already Taken](#when-the-name-is-already-taken)).

## Examples

```bash
# Link the skill in the current folder into Claude Code only
cd skills/loki-query
maia link --harness claude

# Link the current folder into every coding agent (Claude Code / Codex / OpenCode)
maia link

# Link into two specific harnesses at once
maia link --harness claude codex

# Link a skill at a specific path rather than the current folder
maia link ~/skills/loki-query

# Link a skill at a path into Codex only
maia link ~/skills/loki-query --harness codex
```

## Name

A skill is named after its folder.
If the `name:` in the folder's `SKILL.md` doesn't match the folder name, the command stops and exits.
This is the same rule [`maia validate`](../validate/) checks.

## When the Name Is Already Taken

Before linking, the command checks whether a skill of the same name is already there in each harness you're linking into.
There are three cases:

1. **Nothing is there** - no skill of that name is installed, so the link is created.

2. **A managed skill is there** - a skill installed from a marketplace like `ai.integrations`.
   The command stops and points you at `maia uninstall`.
   Uninstall it (or rename your skill folder), then link.
   When you later want the marketplace skill back, re-install it with [`maia install`](../install/).

3. **Anything else is there** - another link of the same name, or an untracked file/folder sitting in the skills directory.
   The command leaves it untouched and tells you to [`maia unlink`](../unlink/) it (if it's a link) or remove it by hand.

## Where Links Are Recorded

Each link is recorded in `~/.maia/maia.json`, so [`maia unlink`](../unlink/) and [`maia status`](../status/) know what you've linked and can restore a skill your link was hiding.
The manifest root defaults to `~/.maia` and can be overridden with the `MAIA_HOME` env var.


# Maia Unlink 
---
title: maia unlink
description: Remove a link you made, and restore any installed skill it was hiding
weight: 80
---

`maia unlink` removes a link you made with [`maia link`](../link/) once you're done developing a skill locally.

Not sure what you've got linked?
Start with [`maia status`](../status/), which lists every link you have - its name, the harness it lives in, and whether it's still working.
Pick the one you're finished with and remove it by name:

```bash
maia unlink <skill-name>
```

`maia unlink` only removes links made by `maia link`.
It never deletes an installed skill or a folder it doesn't know about.

You can remove a link by name, by its source path, or remove all of them with `--all`.
A broken link can still be removed by name.
If you don't pass `--harness`, the link is removed from every harness it's in.

## Usage

```bash
maia unlink [-h] [--harness HARNESS ...] [--all] [name_or_path]
```

## Arguments

| Argument | Notes |
| --- | --- |
| `name_or_path` | The skill to unlink. Accepts either the skill name (its folder name) or the source path you linked from. Omit it only when using `--all`. |

## Options

| Flag | Default | Notes |
| --- | --- | --- |
| `--harness HARNESS ...` | every harness the link is in | One or more coding agents to remove the link from, space-separated. Omit it to remove the link from every harness it lives in. |
| `--all` | off | Remove every link you have, across every harness. Ignores `name_or_path`. |
| `-h`, `--help` | none | Print help and exit. |

### `--harness`

`--harness` narrows the removal to specific coding agent(s).
It is variadic, so you can pass several at once, space-separated.
Omit it and the link is removed from every harness it currently lives in.

| Value | Skills folder |
| --- | --- |
| `claude` | `~/.claude/skills/` |
| `codex` | `~/.agents/skills/` |
| `opencode` | `~/.claude/skills/` (OpenCode reads the same layout as Claude Code) |

## Examples

```bash
# Unlink by name from every harness it's linked in
maia unlink loki-query

# Unlink by source path instead of name
maia unlink ~/skills/loki-query

# Unlink from a specific harness only, leaving other harnesses linked
maia unlink loki-query --harness claude

# Remove every link you have
maia unlink --all
```

## What It Does Not Do

- It does not delete installed skills or any folder it didn't create.
- It does not touch untracked files or folders that happen to share the skill's name - only links `maia link` recorded.


# Maia Status
---
title: maia status
description: See your links, whether they're working, and which git branch is live
weight: 90
---

`maia status` shows the state of your links: whether each one is working and which version of the skill the agent is actually loading.
With no argument, it shows every link you have.
Pass a skill name to look up one link, or a path to check the folder at that location.

It's the companion to [`maia link`](../link/) and [`maia unlink`](../unlink/): link a folder, check on it here, and unlink it when you're done.

## Usage

```bash
maia status [-h] [--json] [name_or_path]
```

## Arguments

| Argument | Default | Notes |
| --- | --- | --- |
| `name_or_path` | all links | A skill name to show one link, or a path to check the folder at that location. Omit it to show every link you have. |

## Options

| Flag | Default | Notes |
| --- | --- | --- |
| `--json` | off | Print machine-readable JSON instead of the table. |
| `-h`, `--help` | none | Print help and exit. |

## Examples

```bash
# Show every link you have
maia status

# Look up one link by skill name
maia status loki-query

# Check the folder at a specific path
maia status skills/loki-query

# Emit JSON for scripting
maia status --json
```

## Output

The command prints a table with these columns:

| Column | Meaning |
| --- | --- |
| `NAME` | The skill's name (its folder name) |
| `HARNESS` | Which harness the link lives in (Claude Code, Codex, OpenCode) |
| `STATE` | Whether the link is working - one of the states below |
| `BRANCH` | The current git branch of the source folder |
| `SOURCE` | Path to the local folder the link points at |

`BRANCH` is worth watching: it shows when a `git checkout` has swapped the skill without you noticing, so you can always tell which version the agent is loading.

## States

A broken link changes the reported `STATE`, never the exit code.
So `maia status` exits `0` whether or not any link is broken - read the `STATE` column (or the JSON) to find problems, rather than the exit code.

| State | Meaning | What to do |
| --- | --- | --- |
| `ok` | The link points to the recorded source and its `SKILL.md` exists | Nothing - the agent is loading your folder |
| `dangling` | The link is correct but the source folder or its `SKILL.md` is gone | Restore the folder, or run [`maia unlink`](../unlink/) to clear it|
| `missing` | The link itself was deleted | Re-create it with [`maia link`](../link/) |
| `replaced` | A real folder or file is now in its place | Remove what's there, then re-link, or rename your skill |
| `foreign` | The link points somewhere other than the recorded source | Re-link from the intended folder, or `maia unlink` it |

## Where It Reads From

`maia status` reads your links from `~/.maia/maia.json` (override the root with `MAIA_HOME`) and then checks each one against what's actually on disk to work out its `STATE` and `BRANCH`.




# Maia Validate

---
title: maia validate
description: Check that a skill is valid before you push it
weight: 60
---

`maia validate` checks that all the contents of a skill are valid.
It checks the `SKILL.md` file and any attachments.

Run it before you push a skill or open a PR: it enforces the same contract MAIA Knowledge expects, so a skill that validates locally won't be rejected for a naming or frontmatter problem later.
It's also the same name-match rule [`maia link`](../link/) applies before it links a folder.

By default it checks the skill in the current folder.
You can also point it at a specific folder, or pass the name of a skill you've [linked](../link/) and it will resolve that to the linked source folder and check it there.
Add `--all` to check every `SKILL.md` under a folder, including skills inside plugins.

## Usage

```bash
maia validate [-h] [--all] [--json] [path]
```

If you don't give a `path`, it uses the current folder.

## Arguments

| Argument | Default | Notes |
| --- | --- | --- |
| `path` | current folder | The skill folder to check, or (with `--all`) a folder to search for skills under. You can also give the name of a [linked](../link/) skill, and it checks that link's source folder. |

## Options

| Flag | Default | Notes |
| --- | --- | --- |
| `--all` | off | Check every `SKILL.md` found under `path`, not just the one folder. |
| `--json` | off | Print results as JSON instead of human-readable text. |
| `-h`, `--help` | none | Print help and exit. |

## Examples

```bash
# Check the skill in the current folder
cd skills/loki-query
maia validate

# Check a skill at a specific path
maia validate skills/loki-query

# Check a skill you've linked, by its name (resolves to the linked source folder)
maia validate loki-query

# Check every skill in a folder (including skills inside plugins)
maia validate --all managed/

# Output as JSON (for CI / scripts)
maia validate --json
```

## What It Checks

There are two kinds of result.
Errors stop the skill from being accepted and exits.
Warnings are shown but don't stop anything.

### Errors

| Rule | Requirement |
| --- | --- |
| `name` length | 64 characters or fewer |
| `name` format | Matches `^[a-z0-9-]+$` (lowercase letters, digits, hyphens) |
| `name` content | Must not contain `anthropic` or `claude` |
| Frontmatter | `SKILL.md` must have YAML frontmatter with both `name:` and `description:` |
| `description` | Non-empty, 1024 characters or fewer, no XML tags (`<...>`) |
| YAML | Frontmatter must be valid YAML that parses to a mapping |
| Encoding | `SKILL.md` must be valid UTF-8 |
| Name match | Frontmatter `name:` must equal the folder name |

The skill is named after its folder, so a `name:` that doesn't match the folder name is rejected.

### Warnings

| Warning | What it flags |
| --- | --- |
| `oversized-file` | Files larger than 5 MB |
| Possible secret | Likely credentials in text files: AWS keys, private-key blocks, GitHub PATs and fine-grained tokens, Slack tokens, and hard-coded `password` / `secret` / `api_key = "..."` assignments |

Binary files (identified by their file extension) are never scanned for secrets.
`node_modules`, `__pycache__`, and hidden folders are skipped.






